% !TeX root = ../main.tex

\section{Band structure of graphene}

\begin{problem}
  In the calculation of the monolayer graphene band structure using
  tight-binding model, if we keep the contribution of the next-nearest
  hopping $t'$.
  \begin{enumext}
    \item Write down the dispersion relation.
    \item Expand the energy spectrum around the Dirac point up to second order.
    Search the literature and check what's trigonal warping effect is in
    a monolayer.
  \end{enumext}
\end{problem}
\begin{solution}\leavevmode
  \begin{enumext}
    \item 
  \end{enumext}
\end{solution}

\begin{problem}
  Graphene has the record room-temperature electron mobility, mainly due to two
  reasons. The first is the high Debye temperature ($\sim\qty{2100}\K$) that
  suppresses phonon scattering arising from the in-plane covalent $sp^2$ bonding
  between carbon atoms. The second one is related to its unique linear Dirac
  band, resulting in massless Dirac fermions ($m^* = 0$) near the Dirac point.
  However, if we apply the definition of the effective
  mass $m^* = \hbar^2\ab(\odv[2]{E(k)}k)^{-1}$ to the band dispersion of
  graphene $(E_\pm(k) = \pm \hbar v_Fk)$, we will get diverging $m^*$ and
  non-differentiable Dirac point. Please explain this contradiction.
\end{problem}
\begin{solution}
  The contradiction arises from applying different definitions of ``mass'' to
  graphene's unique band structure. The description of ``massless Dirac
  fermions'' comes from the analogy to relativistic particles.
  \par \vspace{.5\baselineskip}
  From a massless special relativity particle like a photon to the electrons in
  the graphene, their dispersion relations have the comparasion
  \[
    E = pc \Longrightarrow E = \pm \hbar v_F k = p v_F
  \]
  with $v_F$ playing the role of $c$. Hence, electrons near the Dirac point
  behave as if they have zero rest mass.
  \par \vspace{.5\baselineskip}
  However, the conventional effective mass $m^* = \hbar^2 (\odv[2]E/k)^{-1}$ is
  defined under the parabolic-band approximation $E \propto k^2$. For graphene's
  linear dispersion, $\odv[2]E/k = 0$ away from $k = 0$, leading
  to $m^* \to \infty$ mathematically. At $k = 0$ the dispersion is
  non-differentiable, so the formula breaks down.
  \par \vspace{.5\baselineskip}
  Physically, this infinite effective mass reflects that the electron's speed is
  constant at $v_F$: An external force parallel to its momentum cannot change
  its speed, much like a photon's speed is fixed at $c$. Thus, the ``massless''
  label refers to the relativistic analogy, while the divergent $m^*$ indicates
  the inapplicability of the parabolic effective mass concept to a linear band.
\end{solution}